\documentclass[10pt,a4paper]{amsart}
\usepackage{amsmath,amssymb,amsthm,graphicx}
\usepackage{verbatim}
\newenvironment{reference}{\comment}{\endcomment}
%\newenvironment{reference}{}{}
\newenvironment{slogan}{\comment}{\endcomment}
\newenvironment{history}{\comment}{\endcomment}

% For commutative diagrams we use Xy-pic
\usepackage[all]{xy}

% We use 2cell for 2-commutative diagrams.
\xyoption{2cell}
% Theorem environments.
%
\theoremstyle{plain}
\newtheorem{theorem}[subsection]{Theorem}
\newtheorem{proposition}[subsection]{Proposition}
\newtheorem{lemma}[subsection]{Lemma}

\theoremstyle{definition}
\newtheorem{definition}[subsection]{Definition}
\newtheorem{example}[subsection]{Example}
\newtheorem{exercise}[subsection]{Exercise}
\newtheorem{situation}[subsection]{Situation}

\theoremstyle{remark}
\newtheorem{remark}[subsection]{Remark}
\newtheorem{remarks}[subsection]{Remarks}

\numberwithin{equation}{subsection}

% Macros
%
\def\lim{\mathop{\mathrm{lim}}\nolimits}
\def\colim{\mathop{\mathrm{colim}}\nolimits}
\def\Spec{\mathop{\mathrm{Spec}}}
\def\Hom{\mathop{\mathrm{Hom}}\nolimits}
\def\Ext{\mathop{\mathrm{Ext}}\nolimits}
\def\SheafHom{\mathop{\mathcal{H}\!\mathit{om}}\nolimits}
\def\SheafExt{\mathop{\mathcal{E}\!\mathit{xt}}\nolimits}
\def\Sch{\mathit{Sch}}
\def\Mor{\mathop{\mathrm{Mor}}\nolimits}
\def\Ob{\mathop{\mathrm{Ob}}\nolimits}
\def\Sh{\mathop{\mathit{Sh}}\nolimits}
\def\NL{\mathop{N\!L}\nolimits}
\def\CH{\mathop{\mathrm{CH}}\nolimits}
\def\proetale{{pro\text{-}\acute{e}tale}}
\def\etale{{\acute{e}tale}}
\def\QCoh{\mathit{QCoh}}
\def\Ker{\mathop{\mathrm{Ker}}}
\def\Im{\mathop{\mathrm{Im}}}
\def\Coker{\mathop{\mathrm{Coker}}}
\def\Coim{\mathop{\mathrm{Coim}}}

% Boxtimes
%
\DeclareMathSymbol{\boxtimes}{\mathbin}{AMSa}{"02}

%
% Macros for moduli stacks/spaces
%
\def\QCohstack{\mathcal{QC}\!\mathit{oh}}
\def\Cohstack{\mathcal{C}\!\mathit{oh}}
\def\Spacesstack{\mathcal{S}\!\mathit{paces}}
\def\Quotfunctor{\mathrm{Quot}}
\def\Hilbfunctor{\mathrm{Hilb}}
\def\Curvesstack{\mathcal{C}\!\mathit{urves}}
\def\Polarizedstack{\mathcal{P}\!\mathit{olarized}}
\def\Complexesstack{\mathcal{C}\!\mathit{omplexes}}
% \Pic is the operator that assigns to X its picard group, usage \Pic(X)
% \Picardstack_{X/B} denotes the Picard stack of X over B
% \Picardfunctor_{X/B} denotes the Picard functor of X over B
\def\Pic{\mathop{\mathrm{Pic}}\nolimits}
\def\Picardstack{\mathcal{P}\!\mathit{ic}}
\def\Picardfunctor{\mathrm{Pic}}
\def\Deformationcategory{\mathcal{D}\!\mathit{ef}}
\usepackage{cleveref}
\crefname{originalsection}{original section}{original sections}
\Crefname{originalsection}{Original section}{Original sections}
\makeatletter
\newlabel{example-proper-over-local}{{0AWI}{source}{}{sourcepointer-1}{}}
\newlabel{example-proper-over-local@cref}{{[originalsection][0][]0AWI}{[1][0][]source}{}{}{}}
\makeatother
\crefname{originallemma}{original lemma}{original lemmas}
\Crefname{originallemma}{Original lemma}{Original lemmas}
\makeatletter
\newlabel{lemma-good-dualizing-dimension-function}{{0AWM}{source}{}{sourcepointer-2}{}}
\newlabel{lemma-good-dualizing-dimension-function@cref}{{[originallemma][0][]0AWM}{[1][0][]source}{}{}{}}
\makeatother
\crefname{originallemma}{original lemma}{original lemmas}
\Crefname{originallemma}{Original lemma}{Original lemmas}
\makeatletter
\newlabel{dualizing-lemma-sitting-in-degrees}{{0A7U}{source}{}{sourcepointer-3}{}}
\newlabel{dualizing-lemma-sitting-in-degrees@cref}{{[originallemma][0][]0A7U}{[1][0][]source}{}{}{}}
\makeatother
\crefname{originallemma}{original lemma}{original lemmas}
\Crefname{originallemma}{Original lemma}{Original lemmas}
\makeatletter
\newlabel{dualizing-lemma-depth-dualizing-module}{{0AWE}{source}{}{sourcepointer-4}{}}
\newlabel{dualizing-lemma-depth-dualizing-module@cref}{{[originallemma][0][]0AWE}{[1][0][]source}{}{}{}}
\makeatother
\crefname{originallemma}{original lemma}{original lemmas}
\Crefname{originallemma}{Original lemma}{Original lemmas}
\makeatletter
\newlabel{lemma-dualizing-module}{{0AWK}{source}{}{sourcepointer-5}{}}
\newlabel{lemma-dualizing-module@cref}{{[originallemma][0][]0AWK}{[1][0][]source}{}{}{}}
\makeatother
\title{Vanishing, support and (S\_2) properties of the dualizing module over a local base}
\author{The Stacks Project Authors}
\date{stacks-project:duality.tex:local-snapshot-20261012; GNU FDL 1.2 or later}
\begin{document}
\maketitle
\noindent\textit{Editorial note.} All delivered mathematical text is the author's own, copied from the pinned Stacks Project chapter file duality.tex; no statement or proof step was rewritten, reordered or removed.\par\medskip
\noindent\textit{Editorial note.} Preamble provenance: the pinned chapter file opens with an input of the shared preamble and carries no preamble text of its own, so the project's single Stacks-wide preamble staged at sources/stacks-duality/preamble.tex is excerpted verbatim under the recipe fixed by LANE\_BRIEF 22.1, as three spans, lines 8-13 (verbatim package and the reference/slogan/history environments), lines 14-18 (Xy-pic setup) and lines 154-217 (the native theorem declarations and mathematical macros). Nothing was added, removed or redefined; the wrapper class and the amsmath/amssymb/amsthm/graphicx base packages are supplied by the assembler.\par\medskip
\noindent\textit{Editorial note.} The counted statement and proof contain no cross reference, no citation and no included file, so the delivered body is self contained.\par\medskip
\noindent\textit{Editorial note.} Canonical location of the counted proposition: Stacks tag 0AWN (sources/stacks-duality/pinned-original-source.tex line 6125).\par\medskip
\noindent\textit{Allowed standard prerequisites (editorial).}\begin{itemize}
\item The dimension function of the normalised dualizing complex. Same chapter, Lemma labelled lemma-good-dualizing-dimension-function, cited at proof lines 6156-6160. Identifies the dimension function with the dimension of the closure at closed points.
\item Local vanishing and depth statements for a dualizing complex. Dualizing Complexes, Lemmas labelled dualizing-lemma-sitting-in-degrees and dualizing-lemma-depth-dualizing-module, cited at proof lines 6166-6170. Give the four properties after localising.
\item The dualizing module of a proper scheme over a local base. Same chapter, Lemma labelled lemma-dualizing-module, cited at proof lines 6170-6172. Alternative deduction of the last two statements.
\end{itemize}\medskip
\ifdefined\hypertarget\hypertarget{sourcepointer-1}{}\fi
\noindent\textit{Source reference.} Original section 0AWI: sources/stacks-duality/pinned-original-source.tex, source lines 6070--6070. Same-chapter pointer into the author environment example at source line 6070; invoked at proof line None of the counted proof. It is retained only as a navigation pointer; every load-bearing step is inside the delivered span.\par\medskip
\ifdefined\hypertarget\hypertarget{sourcepointer-2}{}\fi
\noindent\textit{Source reference.} Original lemma 0AWM: sources/stacks-duality/pinned-original-source.tex, source lines 5856--5856. Same-chapter pointer into the author environment lemma at source line 5856; invoked at proof line 6141 of the counted proof. It is retained only as a navigation pointer; every load-bearing step is inside the delivered span.\par\medskip
\ifdefined\hypertarget\hypertarget{sourcepointer-3}{}\fi
\noindent\textit{Source reference.} Original lemma 0A7U: sources/stacks-chow/tags.txt, source lines 12892--12892. Cross-chapter standard prerequisite located in the project tag list staged in the local snapshot; invoked at proof line 6149 of the counted proof. It is retained only as a navigation pointer; every load-bearing step is inside the delivered span.\par\medskip
\ifdefined\hypertarget\hypertarget{sourcepointer-4}{}\fi
\noindent\textit{Source reference.} Original lemma 0AWE: sources/stacks-chow/tags.txt, source lines 13717--13717. Cross-chapter standard prerequisite located in the project tag list staged in the local snapshot; invoked at proof line 6150 of the counted proof. It is retained only as a navigation pointer; every load-bearing step is inside the delivered span.\par\medskip
\ifdefined\hypertarget\hypertarget{sourcepointer-5}{}\fi
\noindent\textit{Source reference.} Original lemma 0AWK: sources/stacks-duality/pinned-original-source.tex, source lines 6101--6101. Same-chapter pointer into the author environment lemma at source line 6101; invoked at proof line 6153 of the counted proof. It is retained only as a navigation pointer; every load-bearing step is inside the delivered span.\par\medskip
\begin{lemma}
\label{lemma-vanishing-good-dualizing}
Let $X/A$ with $\omega_X^\bullet$ and $\omega_X$ be as in
Example \ref{example-proper-over-local}. Then
\begin{enumerate}
\item $H^i(\omega_X^\bullet) \not = 0 \Rightarrow
i \in \{-\dim(X), \ldots, 0\}$,
\item the dimension of the support of $H^i(\omega_X^\bullet)$ is at most $-i$,
\item $\text{Supp}(\omega_X)$ is the union of
the components of dimension $\dim(X)$, and
\item $\omega_X$ has property $(S_2)$.
\end{enumerate}
\end{lemma}

\begin{proof}
Let $\delta_X$ and $\delta_S$ be the dimension functions associated to
$\omega_X^\bullet$ and $\omega_S^\bullet$ as in
Lemma \ref{lemma-good-dualizing-dimension-function}.
As $X$ is proper over $A$, every closed subscheme of $X$ contains
a closed point $x$ which maps to the closed point $s \in S$
and $\delta_X(x) = \delta_S(s) = 0$. Hence
$\delta_X(\xi) = \dim(\overline{\{\xi\}})$ for any point
$\xi \in X$. Hence we can check each of
the statements of the lemma by looking at what happens over
$\Spec(\mathcal{O}_{X, x})$ in which case the result follows
from Dualizing Complexes, Lemmas \ref{dualizing-lemma-sitting-in-degrees} and
\ref{dualizing-lemma-depth-dualizing-module}.
Some details omitted.
The last two statements can also be deduced from
Lemma \ref{lemma-dualizing-module}.
\end{proof}

\end{document}
